\section{Введение}

В настоящее время актуальной задачей для развития сферы телекоммуникаций
является создание группировок спутников для покрытия мобильной связью 
поверхности Земли. Такое покрытие должно существенно упростить обеспечение
связи в труднодоступных местах, тем самым обеспечив дальнейшее развитие науки и
техники.  

Современные технологии космической связи требуют высокой надежности и
эффективности передачи данных, что делает актуальным исследование
различных факторов, влияющих на качество сигналов.

Одним из таких факторов
является эффект Доплера, который возникает в результате относительного
движения источника и приемника сигналов. В современных системах передачи
данных [1], использующих ортогональное частотное деление (OFDM,
Orthogonal Frequency Division Multiplexing), влияние эффекта Доплера может
существенно сказаться на характеристиках принимаемого сигнала, таких как
уровень искажений, скорость передачи данных и устойчивость к помехам.

OFDM является одной из наиболее перспективных технологий для
реализации высокоскоростной передачи информации в условиях ограниченной
полосы пропускания и сложных каналов связи.

Благодаря своей способности
эффективно противостоять межсимвольной интерференции и многолучевому
распространению, OFDM широко применяется в системах связи.

% здесь можно добавить множество ссылок на работы

Однако при
наличии значительных скоростей движения приемников или передатчиков,
таких как спутники или космические аппараты, эффект Доплера приводит к
смещению частот поднесущих, что в свою очередь вызывает проблемы с
синхронизацией и декодированием информации.

Наряду со смещением частот, для борьбы с которым есть множество методов,
возникает эффект масштабирования спектра, который приводит к возникновению
интерференции между поднесущими. Так как основной принцип OFDM состоит в ортогональности
поднесущих, данный метод особенно чувствителен к такой интерференции.

% Начинаем перечислять ссылки
Множество работ было посвящено смещению спектра, например эти (). 

% Тут можно привести описание этих работ

Следующие работы исследовали влияние масштабирования спектра ().

% Описание этих работ

Но в этих работах мастабирование считалось в следствии переотражений сигнала, в следствии рассеяния.
В данной же работе будет рассмотрено масштабирование спектра без переотражений, так как в спутниковой
связи ... 

